\section*{Розділ VI. СР ТПК факультетів та гуртожитків}
\addcontentsline{toc}{section}{Розділ VI. СР ТПК факультетів та гуртожитків}

\subsection*{6.1. СР ТПК факультетів/інститутів}
\addcontentsline{toc}{subsection}{6.1. СР ТПК факультетів/інститутів}
    6.1.1. Студентська рада ТПК факультету/інституту (далі -- СР ТПК факультету) є органом тимчасового студентського самоврядування на рівні факультету/інституту Університету.

    6.1.2. Перелік СР ТПК факультетів та їх голови затверджуються окремим рішенням КСУ.

    6.1.3. Голова СР ТПК факультету призначається та знімається колегіальним рішенням відповідно до п. 4.3 цього Положення за поданням Голови СР ТПК КАІ.

    6.1.4. Головою СР ТПК факультету може бути студент відповідного факультету/інституту Університету денної форми навчання.

    6.1.5. Повноваження Голови СР ТПК факультету припиняються у разі:

        \begin{enumerate}[label=\alph*)]
            \item Зняття з посади відповідно до п. 6.1.3;
            \item Подання особистої заяви про складення повноважень;
            \item Втрати статусу студента відповідного факультету/інституту Університету денної форми навчання;
            \item Припинення діяльності ТПК.
        \end{enumerate}

    6.1.6. Голова СР ТПК факультету:

        \begin{enumerate}[label=\alph*)]
            \item Представляє інтереси студентів відповідного факультету/інституту перед його адміністрацією;
            \item Забезпечує інформування студентів факультету/інституту з питань, що стосуються їхніх прав, обов'язків та інтересів;
            \item Бере участь у роботі колегіальних органів факультету/інституту відповідно до Положення про ОСС;
            \item Підзвітний Голові СР ТПК КАІ та звітує перед ним про свою діяльність;
            \item Здійснює інші повноваження за дорученням Голови СР ТПК КАІ.
        \end{enumerate}

\subsection*{6.2. СР ТПК гуртожитків}
\addcontentsline{toc}{subsection}{6.2. СР ТПК гуртожитків}
    6.2.1. Студентська рада ТПК гуртожитку (далі -- СР ТПК гуртожитку) є органом тимчасового студентського самоврядування на рівні гуртожитку Університету.

    6.2.2. Перелік СР ТПК гуртожитків та їх голови затверджуються окремим рішенням КСУ.

    6.2.3. Голова СР ТПК гуртожитку призначається та знімається колегіальним рішенням відповідно до п. 4.3 цього Положення за поданням Голови СР ТПК СМ.

    6.2.4. Головою СР ТПК гуртожитку може бути студент Університету денної форми навчання, який мешкає у відповідному гуртожитку.

    6.2.5. Повноваження Голови СР ТПК гуртожитку припиняються у разі:

        \begin{enumerate}[label=\alph*)]
            \item Зняття з посади відповідно до п. 6.2.3;
            \item Подання особистої заяви про складення повноважень;
            \item Втрати статусу студента Університету денної форми навчання або припинення проживання у відповідному гуртожитку;
            \item Припинення діяльності ТПК.
        \end{enumerate}

    6.2.6. Голова СР ТПК гуртожитку:

        \begin{enumerate}[label=\alph*)]
            \item Представляє інтереси мешканців відповідного гуртожитку перед адміністрацією студмістечка та гуртожитку;
            \item Забезпечує інформування мешканців гуртожитку з питань, що стосуються їхніх прав, обов'язків та умов проживання;
            \item Сприяє дотриманню мешканцями правил внутрішнього розпорядку гуртожитку;
            \item Підзвітний Голові СР ТПК СМ та звітує перед ним про свою діяльність;
            \item Здійснює інші повноваження за дорученням Голови СР ТПК СМ.
        \end{enumerate}
